\documentclass[12pt]{article}
\usepackage[a3paper, margin=1in]{geometry}
\usepackage{amsmath, amssymb, amsthm, ulem}

\begin{document}

\section*{Domácí úkol 7.1}

Vypracoval: Daniel \textit{``Randál"} Ransdorf \hfill Podpis: \rule{4cm}{0.4pt}

\begin{enumerate}
  
  \item Úloha je opravdu podobná poslední úloze z midtermu. 
    
    Kolik je v $\mathbb{Z}_7^3$ různých vektorů?
      Mám tři pozice, do kterých vybírám ze 7 prvků --- $7^3 = 343$
      
    Posloupnost nesmí obsahovat nulový vektor, jinak by byla závislá.
    Když máme nenulový vektor $\vec{x}$, tak díky prostosti násobení v tělesech máme
      v $\mathbb{Z}_7^3$ dalších 6 násobku toho vektoru.
    Když máme dva nezávislé vektory $\vec{x}_1,\vec{x}_2$ v $\mathbb{Z}_7^3$, pak jejich lineární obal
      (množina všech vektorů závislých na $\vec{x}_1,\vec{x}_2$) má velikost
      $|Z_7^2| = |Z_7|^2 = 49$, protože oba vektory určují jiný ``směr".

    Začněme postupně volit vektory v posloupnosti:

    \begin{enumerate}
      \item $\vec{v}_1$: Zatím máme jen omezení, že nesmíme zvolit nulový vektor.
        Tedy máme $343-1=342$ možností.
      \item $\vec{v}_2$: Nesmíme zvolit vektor závislý na $\vec{v}_1$,
        tedy nesmíme zvolit jeho násobek (-7). Tedy máme $343-7=336$ možností.
      \item $\vec{v}_3$: Nesmíme zvolit vektor závislý na $\vec{v}_1,\vec{v}_2$ (-49).
        Tedy máme $343-49=294$ možností.
    \end{enumerate}

    Postupnou volbou dostáváme,
      že máme $342*336*294 = \underline{\underline{33784128}}$ možných posloupností.
    

\end{enumerate}

\end{document}
